\section{Limitations}
\label{sec:limitations}

\methodname{} depends on the quality of OpenIE extraction and phrase
normalization. Missing, overly generic, or incorrectly normalized
facts can remove useful bridge edges or introduce spurious endpoint
matches, especially for rare entities and relations with non-canonical
surface forms. The query-side anchor and requirement extraction also
uses a local instruction model, so errors in that analysis can make the
query-local graph too narrow or steer coverage reranking toward the
wrong requirements.

The current experiments use English Wikipedia-style corpora and a
fixed top-five reader budget. Results may differ for noisier domains,
longer documents, or settings where the reader can consume much larger
contexts. The method also incurs offline indexing cost from OpenIE
extraction and graph construction, and a deployed system should report
latency and memory under its target corpus scale. Finally, any
evidence-selecting retriever can be misused to assemble selective
evidence bundles, so downstream applications should expose retrieved
sources and preserve uncertainty when answers are presented to users.
